\section{Comparing the Three Interpretations}
\runninghead{Comparing the Three Interpretations}

All three interpretations begin from the same appointed prayers and readings: the pressing danger of our own sins named in the Collect, the call to wake from sleep and put on Christ in Romans~13, the unashamed waiting of Psalm~24, the longing for God's mercy and saving gift in Psalm~84, the sovereign majesty of the Son of Man in Luke~21, and the Trinitarian praise of the Sunday Preface. They diverge not on Christian doctrine, but on the exegetical question each asks of the texts and on how the Fathers resolve specific textual difficulties in Luke~21 and Psalms~24 and~84.

\begin{comparisontable}{Point of Comparison}{Strict Judge}{Incarnational \& Marian}{Conversion \& Altar Bridge}
\textbf{Governing question} & How does cosmic ruin become a command to rejoice, and what fruit does the fig tree put forth? & Why open Advent at St.~Mary Major with a cloud, all the trees, and our earth yielding fruit? & How does grace transform sinners today, between Bethlehem and Judgment, so human soil bears fruit?\\
\textbf{Luke~21:25--33} & Literal cosmic heralds; angelic judges; ironic ``fruit of the world is its ruin'' (Gregory); bodily resurrection (Cyril); renewal \latin{per imaginem} through fire (Bede on v.~33, applying Augustine, \work{De~civ.~Dei}~XX.14, 16, 24). & Symbolic darkening of Christ, Church, and saints; \latin{nube} = prophets and Virgin Mary; ``all the trees'' (\latin{et omnes arbores}) = Israel and all nations blossoming (Ambrose). & Christ's advent \latin{in singulis} shaking the soul's powers by an increase of grace (Ambrose); Advent vigilance against surfeiting and worldly cares (Leo, Luke~21:34).\\
\textbf{Ps.~84:13} & Divine goodness sustaining faithful endurance through the world's twilight. & \latin{Benignitas} = Holy Spirit and untouched virgin \latin{terra nostra} = Mary yielding Christ (Bellarmine, Sicard); Mary ``made fruitful by the dew of heaven'' (Gu\'eranger); Christ as the Fruit of Ps.~84:13 and Ps.~131:11 (Bernard, \work{Missus est}, Hom.~1). & Old Latin \latin{suavitatem} watering hearts ploughed by confession (Augustine); Eucharistic pledge refreshing arid hearts (Schuster); good works (Durandus).\\
\textbf{Ps.~24:1--4} & Tearing the soul from earthly briers (Bellarmine) into watchful, unashamed expectation. & Waiting Israel and the whole Church lifting up her soul for the promised Redeemer. & \latin{Totus Christus} humbled from pride (Augustine); contemplative ascent and commandments (\latin{viae}) vs.\ counsels (\latin{semitae}, Aquinas).\\
\textbf{Rom.~13:11--14a} & Runners straining toward the prize; glittering \latin{arma lucis} making the believer soldier and bridegroom (Chrysostom). & Dawn broken at the Sun's birth; putting on Christ because the Son of God became man so man might become a son of God (Durandus). & Instant personal break with carnal habit to put on Christ (Augustine, \work{Confessions}~VIII.12.29; noting Missal ends at v.~14a).\\
\textbf{Practical emphasis} & Detachment from a passing age, sobriety, and fearless longing for bodily resurrection. & Wonder at the Incarnation and Mary's role, linking Bethlehem directly to the Church's mission. & Confession, humility under prevenient grace, moral growth, and Eucharistic purification at the altar.\\
\end{comparisontable}

The Fathers diverge directly on two passages. On Luke~21:29--30, St.~Gregory reads the fig tree's fruit ironically as the world's decay announcing God's kingdom, whereas St.~Ambrose reads Luke's unique ``and all the trees'' positively as Israel and the nations blossoming in faith toward resurrection (while keeping a warning against barren leaves). On Psalm~84:13, St.~Robert Bellarmine (followed by Sicard of Cremona) expounds \latin{benignitas} of the Holy Spirit and \latin{terra nostra} of the untouched Virgin Mary yielding Christ---with St.~Bernard naming Christ the Fruit of Psalm~84:13 in \work{Missus est} Homily~1 (p.~59) and comparing Mary's womb to an unploughed field in his \work{Second Sermon for Advent} (p.~19)---whereas St.~Augustine, working from his Old Latin text \latin{suavitatem}, expounds the verse of sinful hearts ploughed by confession and watered by interior delight in righteousness, and Bl.~Ildefonso Schuster applies it directly to Holy Communion refreshing arid souls.

St.~Bernard (\work{In~Adventu}~4.1, pp.~32--33) urges Christians to keep both visible comings together in meditation, pondering ``how much we have received by the first, how much we are promised at the second,'' comparing them to the Bridegroom's two arms. Taken together, the three interpretations correspond to the threefold coming of Christ (\latin{triplex adventus}) summarized by John Beleth (\work{Summa}~LXIV, PL~202, col.~72) and preached by St.~Bernard (\work{In~Adventu}~3.3--4, 5.1): Christ came in human flesh through the Virgin Mary (the second interpretation), comes daily within human hearts by the Holy Spirit and grace---joined in the third interpretation with Schuster's Eucharistic exposition---and will come openly in majesty to judge the world and renew creation (the first).
